\section{The exact real-period conjugate-pair collision and its two inverse spectra}
This calculation retains the original unit, quartet, periods and all
Gamma metrics of PCL1--10, OFM1--17, WCF1--6 and GFC1--21
\cite{PCLOriginal,NC,Cumulative,WCF}. It calculates the exact local
matrix on the locus of simple conjugate-pair collisions defined below.
The RPZ full-series contraction proof supplies an actual positive
real-period point on this locus, and RPE1--26 specializes every
constant and four-label map there. The general calculation below
also determines the complete exceptional coefficient locus.

\subsection{Joint conjugation in both original variables}
Retain $a=\alpha+i\chi\ne0$ with real $\alpha,\chi$, the unchanged
$U=\operatorname{diag}(a,\bar a,\bar a,a)$, and the original
ordered quartet in $V$. Let $J$ exchange primary indices $1,2$ and
$3,4$. For $z=1/u$ and the exponential-curve coordinate $\xi$ put
\[
 L(z)=U^{-1}V\mathcal R(z)^{-1},\qquad
 M_j(z)=L(z)D_5^{-j}L(z)^{-1},\qquad
 f_j(\xi,z)=Q_0(M_j(z)t(\xi)).
 \tag{RPCJ1}
\]
Every PCL2 coefficient is real at the original real $\delta,\gamma$,
so $\overline{\mathcal R(\bar z)}=\mathcal R(z)$.
The original labels and unit give
$\bar V=JV$, $\bar U=JUJ$, $\overline{t(\bar\xi)}=Jt(\xi)$,
and $Q_0(Jv)=-Q_0(v)$. Therefore
\[
 \overline{L(\bar z)}=JL(z),\qquad
 \overline{M_j(\bar z)}=JM_{5-j}(z)J,
 \qquad
 \boxed{\ f_{5-j}(\xi,z)=-\overline{f_j(\bar\xi,\bar z)}.\ }
 \tag{RPCJ2}
\]
To verify the first equality multiply
$\bar U^{-1}\bar V=(JU^{-1}J)(JV)=JU^{-1}V$.
Conjugating the diagonal $D_5^{-j}$ gives $D_5^j$ and proves the
second equality. Substituting the conjugated original $t$ and then
using the exact sign in $Q_0(Jv)$ proves the final equality.
This is the two-variable form of the earlier OFM8--10 identity;
it keeps the original unit fixed rather than replacing its phase.

For $h^\#(\xi,z)=\overline{h(\bar\xi,\bar z)}$, set
$G=f_1f_2$. The two minus signs in RPCJ2 cancel, giving
\[
 E_A(\xi,z)=G(\xi,z)G^\#(\xi,z),\qquad
 E_A(\xi,z)=|G(\xi,z)|^2\ge0
                 \quad(\xi,z\in\mathbb R).
 \tag{RPCJ3}
\]
At every real period, the full moment formula is
\[
 \mu_n(z)=\sum_{r=0}^n\binom nr
        \partial_\xi^rG(0,z)\,
                  \overline{\partial_\xi^{n-r}G(0,z)}.
 \tag{RPCJ4}
\]
Each derivative of $G$ is itself the complete original factor sum
$\partial_\xi^rG=\sum_{h=0}^r\binom rh
(\partial_\xi^hf_1)(\partial_\xi^{r-h}f_2)$.
Thus RPCJ4 retains every original moment and repeated contribution.
The finite nine-frequency test and order bound from PZ12 and OFM2--5
apply with their existing provenance
\cite{PeriodZeroOriginal,Cumulative}: when the symbol is nonzero,
$n_4=n_1$, $n_3=n_2$, and its order is
$v=2(n_1+n_2)\le32$ with positive leading moment.
An identically zero factor instead gives an identically zero symbol,
as in GFC13--14; no finite order is assigned in that case.

\subsection{The exact simple-pair locus and every moment there}
Define the local locus under study by an original real
$z_0\ne0$ at which exactly one of $f_1(0,z_0)$ and $f_2(0,z_0)$
vanishes, the other is nonzero, and the vanishing factor has nonzero
first derivatives with respect to both $z$ and $\xi$ at $(0,z_0)$.
RPCJ2 forces its conjugate partner to vanish with the same simple
orders, while the other partner is nonzero. This definition retains
all four factors; it is not a certificate that this locus is nonempty.
Write the complete actual jets
\[
 \begin{gathered}
 A=G_z(0,z_0)\ne0,\quad B=G_\xi(0,z_0)\ne0,\\
 C=G_{zz}(0,z_0),\quad H=G_{\xi z}(0,z_0),\quad
 F=G_{\xi\xi}(0,z_0),\qquad \tau=z-z_0.
 \end{gathered}\tag{RPCJ5}
\]
For example when $f_1(0,z_0)=0$ and $b=f_2(0,z_0)\ne0$,
these are the full expressions
\[
 \begin{split}
 A&=b f_{1,z},\quad B=b f_{1,\xi},\\
 C&=b f_{1,zz}+2f_{1,z}f_{2,z},\\
 H&=b f_{1,\xi z}+f_{1,z}f_{2,\xi}+f_{1,\xi}f_{2,z},\\
 F&=b f_{1,\xi\xi}+2f_{1,\xi}f_{2,\xi}.
 \end{split}\tag{RPCJ6}
\]
All derivatives on the right are evaluated at $(0,z_0)$ in the
original functions RPCJ1. When $f_2$ is the vanishing factor the
same product derivatives exchange the two displayed labels.

The complete mixed moment formula, valid for every $m,n\ge0$ at
this real period, is
\[
 \partial_z^m\mu_n(z_0)=
 \sum_{r=0}^n\sum_{h=0}^m\binom nr\binom mh
 \partial_\xi^r\partial_z^hG(0,z_0)
 \overline{\partial_\xi^{n-r}\partial_z^{m-h}G(0,z_0)}.
 \tag{RPCJ7}
\]
It follows by differentiating the exact analytic identity RPCJ3;
complex conjugation of the coefficients is justified by the reality
of $z_0$, not by applying conjugation to a complex independent variable.
The leading cases are
\[
 \begin{gathered}
 \mu_0(z)=|A|^2\tau^2+\operatorname{Re}(C\bar A)\tau^3+O(\tau^4),\\
 \mu_1(z)=2\operatorname{Re}(A\bar B)\tau
       +[\operatorname{Re}(B\bar C)+2\operatorname{Re}(H\bar A)]
                                      \tau^2+O(\tau^3),\\
 \mu_2(z_0)=2|B|^2>0,\quad
 \mu_2'(z_0)=2\operatorname{Re}(F\bar A)+4\operatorname{Re}(H\bar B),
 \quad \mu_3(z_0)=6\operatorname{Re}(F\bar B).
 \end{gathered}\tag{RPCJ8}
\]
In particular the actual period-zero multiplicity is $\nu=2$,
while the actual exponential-symbol order is $v=2$.
Both are calculated from nonzero original derivatives, not identified
by their names. The higher moments are all specified by RPCJ7.

\subsection{The exact four-dimensional weighted inverse}
Now use the unchanged old cutoff $D=3$, with $s\in\{1,a_{\rm amp}\}$,
the original centers $a_{\rm amp}/2$ and $a_{\rm amp}/2-4$,
the complete mass $M_s=(2\pi)^{s/2}$, and
$\rho_m=\sqrt{m!/(s/2)_m}$, including $\rho_0=1$.
These are precisely the physical norms of GFC4; the following matrix
is in their original orthonormal frames.
To avoid confusing $G=f_1f_2$ with the conductor series, retain
\[
 \begin{gathered}
 g(t,z)=e^{-4\omega(t)}E_A(\omega(t),z),\qquad
 \omega(t)=-i\arctan t,\\
 B_s(z)=\operatorname{diag}(\rho_0,\ldots,\rho_3)^{-1}
           T_3(g)\operatorname{diag}(\rho_0,\ldots,\rho_3).
 \end{gathered}
 \tag{RPCJ9}
\]
Its four coefficients, with every center correction, are
\[
 \begin{split}
 g_0&=\mu_0,\qquad g_1=-i(\mu_1-4\mu_0),\\
 g_2&=-\mu_2/2+4\mu_1-8\mu_0,\\
 g_3&=i(\mu_3/6-2\mu_2+25\mu_1/3-12\mu_0).
 \end{split}\tag{RPCJ10}
\]
Expanding $\omega=-it+it^3/3+O(t^5)$ and the exact exponential
proves these four coefficients directly.
Define the following actual constants, including their signs:
\[
 \begin{gathered}
 c=|A|^2>0,\quad R=2\operatorname{Re}(A\bar B),\quad
 d=-iR,\quad e=-|B|^2<0,\\
 c_3=\operatorname{Re}(C\bar A),\quad
 d_2=-i[\operatorname{Re}(B\bar C)+2\operatorname{Re}(H\bar A)-4c],\\
 e_1=-\operatorname{Re}(F\bar A)-2\operatorname{Re}(H\bar B)+4R,
 \quad f=i[\operatorname{Re}(F\bar B)-4|B|^2].
 \end{gathered}\tag{RPCJ11}
\]
Equations RPCJ8--10 give
\[
 g_0=c\tau^2+c_3\tau^3+O(\tau^4),\quad
 g_1=d\tau+d_2\tau^2+O(\tau^3),\quad
 g_2=e+e_1\tau+O(\tau^2),\quad g_3=f+O(\tau).
 \tag{RPCJ12}
\]
At $\tau=0$ the nonzero matrix block is
\[
 C_s=\begin{pmatrix}
 e\rho_2/\rho_0&f\rho_3/\rho_0\\
 0&e\rho_3/\rho_1
 \end{pmatrix}.
 \tag{RPCJ13}
\]
Its columns are source indices $2,3$ and its rows are target indices
$0,1$ in the displayed coordinate matrix; equivalently the original
map sends source degrees $2,3$ to target degrees $0,1$.
The matrix $C_s$ is invertible because $e\ne0$.
Thus $B_s(z_0)$ has rank two, and its two positive forward singular
values are exactly the two singular values of this complete block.

The inverse series through degree three is exactly
\[
 \begin{gathered}
 h_0=\frac1{g_0},\quad h_1=-\frac{g_1}{g_0^2},\\
 h_2=\frac{g_1^2-g_0g_2}{g_0^3},\quad
 h_3=\frac{-g_1^3+2g_0g_1g_2-g_0^2g_3}{g_0^4},\\
 B_s^{-1}=\operatorname{diag}(\rho)^{-1}T_3(h)
                               \operatorname{diag}(\rho).
 \end{gathered}
 \tag{RPCJ14}
\]
Multiplication of the two truncated series proves all entries.
Set
\[
       \Delta=d(d^2-2ce),\qquad L=\frac{d^2-ce}{c^3}.
 \tag{RPCJ15}
\]
If $\Delta\ne0$, the only inverse entry with order five is
$h_3\rho_3/\rho_0$. Its leading coefficient is
$-\Delta\rho_3/(c^4\rho_0)$.
Every other entry has order at most four by RPCJ12--14.
It follows that the largest inverse singular value has the exact
limit
\[
 \lim_{z\to z_0}|z-z_0|^5\|B_s(z)^{-1}\|
               =\frac{|\Delta|\rho_3}{c^4\rho_0}>0.
 \tag{RPCJ16}
\]

The exceptional coefficient locus $\Delta=0$ is also completely
determined. Put
\[\theta=\operatorname{Re}(A\bar B)/(|A||B|)\in[-1,1].\]
This quotient is an exact auxiliary value of the two original
derivatives; their magnitudes remain in RPCJ11 and every constant.
Then $\Delta=0$ exactly when
\[
        \theta=0\quad\hbox{or}\quad\theta^2=1/2.
 \tag{RPCJ17}
\]
Indeed $d=-2i|A||B|\theta$ and $ce=-|A|^2|B|^2$.
On this locus $d^2$ equals either $0$ or $2ce$, so
$d^2-ce$ equals $-ce$ or $ce$ and is nonzero.
Thus $L\ne0$ in all exceptional cases.
The complete next coefficient of the top-right inverse entry is
\[
 K=\frac{2cde_1+2c_3de+(2ce-3d^2)d_2-c^2f}{c^4}.
 \tag{RPCJ18}
\]
To prove it, the coefficient of $\tau^4$ in the numerator of
$h_3$ in RPCJ14 is exactly the displayed numerator, whereas its
$\tau^3$ coefficient is $-\Delta=0$. The denominator has leading
coefficient $c^4\tau^8$, so no denominator correction contributes
at order $\tau^{-4}$. Consequently the full leading inverse block is
\[
 N_s=\begin{pmatrix}
 L\rho_2/\rho_0&K\rho_3/\rho_0\\
 0&L\rho_3/\rho_1
 \end{pmatrix},\qquad
 \lim_{z\to z_0}|z-z_0|^4\|B_s(z)^{-1}\|=\|N_s\|>0
                   \quad(\Delta=0).
 \tag{RPCJ19}
\]
For this limit, multiplication by $\tau^4$ makes every other inverse
entry tend to zero, while these four block entries have the stated
limits. Continuity of the operator norm proves RPCJ19 for complex
as well as real approaches to the original real period.

\subsection{Every exterior rank and every inverse singular value}
Independently of $\Delta$, the exact complementary-exterior identity
GFC11 and $\det B_s=\mu_0^4$ prove
\[
 \begin{split}
 \lim_{z\to z_0}|z-z_0|^8\|\bigwedge^2B_s(z)^{-1}\|
   &=\frac{|e|^2\rho_2\rho_3}{c^4\rho_0\rho_1},\\
 \lim_{z\to z_0}|z-z_0|^8\|\bigwedge^3B_s(z)^{-1}\|
   &=\frac{\|C_s\|}{c^4},\\
 \lim_{z\to z_0}|z-z_0|^8\|\bigwedge^4B_s(z)^{-1}\|
   &=c^{-4}.
 \end{split}\tag{RPCJ20}
\]
The first numerator is $|\det C_s|$, because the only nonzero
second exterior singular value of the rank-two limiting matrix is
the product of its two positive singular values. The next numerator
is its largest singular value $\|C_s\|$.
All displayed constants are strictly positive.

For an explicit evaluation of either block norm, if
$T=\left(\begin{smallmatrix}u&v\\0&w\end{smallmatrix}\right)$,
put $S_T=|u|^2+|v|^2+|w|^2$ and $P_T=|uw|^2$.
The two singular values are
\[
 \sigma_\pm(T)=
 \left(\frac{S_T\pm\sqrt{S_T^2-4P_T}}2\right)^{1/2},
 \qquad\|T\|=\sigma_+(T).
 \tag{RPCJ21}
\]
This follows from the trace and determinant of $T^*T$ and keeps
both exact off-diagonal coefficients in RPCJ13 and RPCJ19.

Let the inverse singular values be in descending order.
For $\Delta\ne0$, their pole exponents and leading constants are
\[
 \begin{gathered}
       (5,3,0,0),\\
 \left(
 \frac{|\Delta|\rho_3}{c^4\rho_0},\quad
 \frac{|e|^2\rho_2}{|\Delta|\rho_1},\quad
 \frac1{\sigma_-(C_s)},\quad\frac1{\sigma_+(C_s)}
 \right).
 \end{gathered}\tag{RPCJ22}
\]
The first two constants follow by dividing the first line of RPCJ20
by RPCJ16. The last two follow from the convergence of the two
nonzero forward singular values to $\sigma_\pm(C_s)$.
For $\Delta=0$, the complete answer is instead
\[
       (4,4,0,0),\qquad
 \left(\sigma_+(N_s),\sigma_-(N_s),
       \frac1{\sigma_-(C_s)},\frac1{\sigma_+(C_s)}\right).
 \tag{RPCJ23}
\]
The limiting rank-two inverse block in RPCJ19 is invertible, so both
first constants are nonzero. As a check retaining all weights,
on this locus $|L|=|e|/c^2$ and hence
$|\det N_s|=|\det C_s|/c^4$, exactly RPCJ20's first coefficient.
Thus the exterior exponents are $(5,8,8,8)$ or $(4,8,8,8)$,
respectively. Nonzero first $z$ and $\xi$ derivatives alone do not
select between these two calculated possibilities.

At the collision the order-adjusted map is the exact $v=2$ map
$\mathcal P_5/\mathcal P_1\to\mathcal P_3$ with its original
Gamma quotient norm and diagonal
$(-i)^2\mu_2\rho_{j+2}/(2\rho_j)$, $0\le j\le3$.
It is invertible because $\mu_2=2|B|^2>0$.
This distinct domain is retained alongside the singular fixed old
map; no quotient change is used to erase the fixed-space poles.

Finally $u_0=1/z_0$ is real and nonzero, and
$z-z_0=-(u-u_0)/(uu_0)$ exactly.
Each singular or exterior pole of exponent $p$ above therefore has
the same exponent in the original $u$ coordinate and its displayed
constant multiplied by $|u_0|^{2p}$. This specifies every original
coordinate constant without selecting a candidate real period.
